\section{Design and implementation}
\label{sec:design}
\subsection{Insertion and removal boundaries}
The resident monitor is prepared in Windows~1 before VMX entry. Admission checks
the exposed Intel VMX/EPT capabilities and lifecycle state. In particular,
existing VMX ownership cannot be resolved by changing a visibility bit: the
CR4.VMXE ownership gate remains in place. A guest with a previously active VMM
is outside the insertion contract.

Preparation allocates per-CPU resources and EPT structures and runs VMX
self-tests. The EPTP-switch backend has a separate armed state; generic resource
readiness is insufficient evidence that the intended page-view backend can
operate. The measurement harness checks that specific state before proceeding.
Cold preparation and the resident-start control body have separate timing
boundaries.

A cross-CPU rendezvous captures processor state, programs the VMCS, and returns
through an assembly continuation that resumes the running Windows context as a
guest of \ks{}. The response reports the resident processor count and per-CPU
state. The implementation also coordinates power transitions, topology changes,
and driver unloading with the resident lifecycle. These guards are prerequisites
for keeping the monitor's code and stacks valid; they are not evidence that
every sleep, hot-add, or failure path has been exercised in this evaluation.

Removal of a Windows-only resident instance coordinates VMX exit on all active
processors before releasing resident resources. With descendants present, the
tested cleanup order keeps the monitor running while closing VMware, retires
and drains page overrides, and then stops residency. It does not stop the
monitor underneath live descendant VMX state. An earlier opposite-order cleanup
that timed out is retained as a failed trial.

\subsection{Nested execution and page composition}
The nested VMX layer maintains the intermediate VMM's virtual VMCS state,
constructs the execution VMCS for the descendant, and reflects exits that
belong to the intermediate VMM. Translation composition follows established
nested-virtualization architecture~\cite{turtles}. The page override is applied
when a descendant translation is installed in \ks{}'s composed EPT.

Let $g$ be a descendant guest-physical address, $T(g)$ its source translation
into Windows~1 physical memory, and $E$ the containing monitor's translation.
For a selected aligned page $p$, replacement frame $r$, and currently admitted
lease $\ell$, the intended address component of the composed translation is
\begin{equation}
  C_{\ell}(g)=
  \begin{cases}
    E(r + (g \bmod 4096)), & \lfloor g/4096\rfloor=p/4096
      \text{ and } \ell \text{ authorizes the override},\\
    E(T(g)), & \text{otherwise}.
  \end{cases}
  \label{eq:compose}
\end{equation}
This notation omits permissions and memory-type checks for readability and
does not describe a new hardware translation level. Outer Hyper-V still
virtualizes execution and the physical memory visible to Windows~1. The
override changes the descendant's selected mapping; it does not overwrite
the source data page. Reverting to the source therefore does not require
copying saved data back from the replacement.

The current control interface admits one active lease and retains at most one
retired lease pending reclamation. A new map is refused while either slot is
occupied. This bounded design makes the observed allocation ledger simple; it
is not a many-page transactional memory facility.

\subsection{A lease over process and translation identity}
\label{sec:lease}
A GPA alone is insufficient identity. The caller supplies a current EPT12 root
and GPA and keeps the ordinary write-back RAM page reserved or pinned. Admission
captures the source walk twice, checks agreement, allocates private backing,
and binds the lease to a referenced VMM process object and its creation time.
Holding the object reference prevents a new process with a recycled numeric
PID from reviving the lease.

The version~3 response exposes the original Windows~1 backing and the captured
source path. The walker records the EPTP configuration, addresses and normalized
values of up to four entries, effective permissions, and the resolved frame.
It handles 4\,KiB, 2\,MiB, and 1\,GiB source leaves. When EPT A/D is enabled,
normalization ignores accessed bits and leaf dirty bits maintained by hardware;
it does not ignore a changed frame, permission, page size, memory type, or parent
path. With A/D disabled, those bits are not silently treated as updates.

Matching leases are checked before nested entry and before a shadow-EPT fill,
using per-CPU physical windows. The first observed semantic drift, unreadable
source path, or owner exit makes revocation sticky and advances the policy
generation. A subsequent observation cannot automatically revive the old lease.
The caller must complete retirement and issue a new map.

These checks are sampled. An ABA change that returns every captured bit to its
original value before observation is invisible. A guest reboot or snapshot
restore with the identical translation can also be invisible, as can reuse of
an application object at the same GPA. Callers must remove leases before these
operations. The walker is not a complete architectural validator for every
reserved-bit combination, and the lease is not a security boundary against a
malicious intermediate VMM.

\subsection{Publication, invalidation, and reclamation}
\label{sec:retire}
Figure~\ref{fig:lease} summarizes the control lifecycle. Mutations are serialized.
A fully initialized immutable rule is published through an ordered pointer
exchange, followed by a policy-generation change and an all-CPU invalidation
operation. A successful map response requires the commit path to complete.
Failure after publication takes the rollback/retirement path rather than
freeing storage that a CPU might still reference.

\begin{figure}[tbp]
\centering
\begin{tikzpicture}[font=\small,
  box/.style={draw=blue!45!black,fill=blue!3,rounded corners=2pt,
              align=center,text width=3.7cm,minimum height=19mm,inner sep=5pt},
  arrow/.style={-{Latex[length=2mm]},thick,draw=blue!50!black}]
\node[box] (capture) {Capture twice\par Bind owner + allocate};
\node[box,right=9mm of capture] (publish) {Publish lease\par Commit invalidation};
\node[box,right=9mm of publish] (active) {Active lease\par Sample source path};
\node[box,below=11mm of active] (retire) {Unpublish / revoke\par Retain backing};
\node[box,left=9mm of retire] (drain) {All-CPU drain\par Failure: keep backing\par and retry};
\node[box,left=9mm of drain,fill=green!4,draw=green!40!black] (free) {Drain succeeds\par Reclaim rule + backing};
\draw[arrow] (capture) -- (publish);
\draw[arrow] (publish) -- (active);
\draw[arrow] (active) -- node[right,font=\footnotesize,align=left] {remove /\\observed fault} (retire);
\draw[arrow] (publish.south) -- ++(0,-0.45) -| node[near start,above,font=\footnotesize] {commit failure} (retire.north);
\draw[arrow] (retire) -- (drain);
\draw[arrow] (drain) -- (free);
\end{tikzpicture}

\caption{Page-control lifecycle. Publication, invalidation completion, and
backing reclamation are distinct. A failed drain retains the rule and backing
for retry. Asynchronous owner/path revocation alone is not a completed global
retirement boundary.}
\label{fig:lease}
\end{figure}

Remove first unpublishes the override or handles its revoked state. Per-CPU
composed translations must then be invalidated before the rule, replacement
page, and process reference can be reclaimed. If a drain fails, the retired
slot and backing remain present. Retrying removal repeats the required drain.
Thus, \code{active=0} alone does not prove safe reclamation; \code{retired=0}
and a successful synchronous drain are also required.

The design supports three auditable obligations under the assumptions of
Section~\ref{sec:background}: (1) new shadow fills use only an admitted,
non-revoked lease; (2) a published replacement remains allocated until every
relevant processor completes retirement, or until the stopped lifecycle is
established; and (3) failure is represented as retained or rolled-back state,
not hidden by a successful response. The experiment checks traces and ledger
consequences of these obligations. It is not a formal proof over every possible
interleaving or hardware failure.

Different vCPUs need not observe a policy change simultaneously. A CPU may
retain an old cached translation until its invalidation point, even after
another CPU detects drift. Readers spanning the transition are not guaranteed
an atomic 4\,KiB application update. The backing-retention rule protects the
lifetime of potentially cached mappings; it does not by itself serialize all
application accesses.

\subsection{Implementation and instrumentation}
The monitor is part of a Windows x64 driver built with MSVC/WDK. The nested
implementation, resident lifecycle, EPT composition, page controller, and
portable lease walker are separate modules. The page/metrics response ABI is
versioned; the current lease response is version~3. The main program's embedded
HVM page and command-line tools use a shared command catalog and execution
engine. GUI runtime behavior is not part of the present evaluation.

Measurements expose QueryPerformanceCounter (QPC) boundaries, per-CPU resident
identity, nested exit classes, EPT violations, actual INVEPT-wrapper calls,
shadow-cache events, and rule/replacement allocation counts. Detailed page
events identify capture, publication, drain, rollback, and reclamation. These
counts belong to \ks{}; they are not the total exits of outer Hyper-V or a
complete accounting of every driver pool allocation.

The historical optimization cohort replaces repeated EPT A/D lookups with a
shadow-leaf-indexed ledger, accepts only legal A/D $0\to1$ changes in table
snapshots, enlarges the tracked parent-table capacity from 32 to 160 pages per
CPU, uses CPU-local single-writer counters, avoids duplicate VMREADs, and caches
CPUID leaf~0 per CPU. The current cohort additionally omits two unused telemetry
VMREADs on ordinary CPUID exits. Entry failures and other exits retain the full
snapshot. We evaluate that small change against a same-binary reference flag;
we do not treat the historical bundle as a set of separately isolated causes.
